\chapter{Worked Controller and Verification Cases}
\label{app:worked-cases}

This appendix applies the registered equations to concrete boundary and failure
cases. The values are deterministic worked examples, not additional experimental
measurements.

\section{Threshold Equality and State Encoding}

Consider the controller input
\begin{equation}
(x_K,x_Q,x_L,x_I)=(0.25,0.93,0.75,0.125).
\end{equation}
Round-half-up UNORM16 conversion produces codes
$(16384,60948,49151,8192)$. Every value is exactly on a registered threshold, so
the left-closed boundary rule gives bins
\begin{equation}
(b_K,b_Q,b_L,b_I)=(1,2,3,1).
\end{equation}
The radix-4 address is
\begin{align}
z&=(((1\times4)+2)\times4+3)\times4+1\\
 &=109,
\end{align}
whose binary representation is \texttt{01 10 11 01}. This one example checks the
feature ordering, equality behavior, radix construction, and bit slicing.

\begin{table}[H]
\centering
\caption{Additional state examples and first-match H3 results. Values are shown as
bins because quantization has already been applied.}
\label{tab:worked-h3-states}
\begin{tabularx}{\textwidth}{@{}ccccr c X@{}}
\toprule
$\boldsymbol b_K$ & $\boldsymbol b_Q$ & $\boldsymbol b_L$ &
$\boldsymbol b_I$ & $\boldsymbol z$ & \textbf{H3 action} & \textbf{First matching rule}\\
\midrule
0 & 0 & 3 & 3 & 15 & $\Pi_1$ & Critical occupancy (priority 1)\\
2 & 1 & 3 & 3 & 175 & $\Pi_2$ & Low route-quality bin (priority 2)\\
2 & 2 & 1 & 2 & 166 & $\Pi_1$ & High imbalance (priority 3)\\
2 & 2 & 3 & 1 & 173 & $\Pi_3$ & Saturated load with healthy occupancy/quality
and low imbalance (priority 4)\\
1 & 2 & 3 & 1 & 109 & $\Pi_0$ & No preceding rule matches (default)\\
\bottomrule
\end{tabularx}
\end{table}

The first row also demonstrates priority: even though other stress conditions are
present, $b_K=0$ selects $\Pi_1$. The fourth row shows that the low-latency profile
is requested only when all four conjunctive conditions hold. Here ``low latency''
is the registered profile label for path-cost emphasis, not a promise about the
number of FPGA kernel cycles.

\section{UNORM16 Boundary Cases}

The converter clips before scaling. Representative values are
\begin{table}[H]
\centering
\caption{Round-half-up and saturation examples for UNORM16.}
\label{tab:worked-unorm16}
\begin{tabular}{@{}r r l@{}}
\toprule
\textbf{Real input} & \textbf{Code} & \textbf{Reason}\\
\midrule
$-0.10$ & 0 & Lower saturation\\
0 & 0 & Exact lower endpoint\\
0.125 & 8,192 & $\lfloor0.125(65535)+0.5\rfloor$\\
0.25 & 16,384 & First occupancy/load threshold\\
0.50 & 32,768 & Middle occupancy/load threshold\\
0.75 & 49,151 & Third occupancy/load threshold\\
0.90 & 58,982 & First route-quality threshold\\
0.93 & 60,948 & Second route-quality threshold\\
0.96 & 62,914 & Third route-quality threshold\\
1.00 & 65,535 & Exact upper endpoint\\
1.20 & 65,535 & Upper saturation\\
\bottomrule
\end{tabular}
\end{table}

Comparing integer codes rather than reconstructed real numbers removes ambiguity at
thresholds. For example, a code of 60,948 enters route-quality bin 2, whereas 60,947
remains in bin 1.

\section{Reward Examples}

Suppose a successful route has bounded quality $b=0.96$, post-decision imbalance
$I^{\rm post}=0.25$, and three hops. Then
\begin{equation}
U_Q=\frac{0.96-0.90}{0.10}=0.6,
\qquad U_B=1-0.25=0.75.
\end{equation}
With no profile change, the evaluation reward is
\begin{equation}
r=4+2(0.6)+0.75-0.25(3)=5.20.
\end{equation}
If the profile changes on the same decision, the switch penalty gives $r=4.95$.
If no candidate is feasible, the route terms are not evaluated: the reward is
$-4.00$ without a switch and $-4.25$ with one.

For the controlled training target, only the balance coefficient changes from one
to four. The same successful non-switching transition therefore has target reward
\begin{equation}
r_{\rm train}=4+2(0.6)+4(0.75)-0.25(3)=7.45.
\end{equation}
The final validation table does not report this 7.45-scale quantity; it recomputes
5.20 on the original evaluation scale. This worked case illustrates why training
targets and reported evaluation utilities must be named separately.

\begin{table}[H]
\centering
\caption{Direction of each term in the worked reward.}
\label{tab:worked-reward-directions}
\begin{tabularx}{\textwidth}{@{}p{0.25\textwidth}p{0.21\textwidth}X@{}}
\toprule
\textbf{Change} & \textbf{Reward effect} & \textbf{Scope}\\
\midrule
Successful rather than blocked & $+8$ before route utilities/costs & Service outcome\\
Higher quality above 0.90 & Positive, clipped at $+2$ & Abstract experimental code\\
Smaller post-decision imbalance & Positive & Synthetic balance objective\\
One extra hop & $-0.25$ & Topological path length\\
Profile change & $-0.25$ & Controller stability\\
\bottomrule
\end{tabularx}
\end{table}

\section{Candidate Normalization Edge Cases}

For two feasible candidates, the range denominator is one greater than the raw
difference. If cost codes are $98{,}304$ and $131{,}072$, the larger cost receives
\begin{equation}
\operatorname{round}_{\rm half\ up}
\left(\frac{32768\times65535}{32769}\right)=65533,
\end{equation}
not 65,535. If both cost codes are equal, both normalized cost penalties are zero
and the cost objective cannot distinguish the routes. The comparator then depends
on the remaining profile-weighted objectives and, if the complete scores tie, the
raw cost/hop/slot ordering.

With only one feasible candidate, every objective range is constant, all three
normalized penalties are zero, and the candidate wins regardless of profile. This
does not mean the route has perfect physical quality; it means there is no feasible
alternative in the supplied set against which to normalize it. The short cycle
cases in H1--H4 occur when such constant ranges bypass iterative division.

\begin{table}[H]
\centering
\caption{Deterministic winner behavior for common edge cases.}
\label{tab:worked-winner-cases}
\begin{tabularx}{\textwidth}{@{}p{0.28\textwidth}p{0.21\textwidth}X@{}}
\toprule
\textbf{Candidate condition} & \textbf{Status} & \textbf{Selection behavior}\\
\midrule
No feasible candidate & \texttt{NO\_PATH} & Path 255; score/quality zero\\
Exactly one feasible candidate & \texttt{OK} & Select it; normalized ranges are zero\\
Two candidates, unequal scores & \texttt{OK} & Select smaller 18-bit score\\
Equal scores, unequal raw cost & \texttt{OK} & Select smaller UQ16.16 path cost\\
Equal score/cost, unequal hops & \texttt{OK} & Select smaller hop count\\
Equal score/cost/hops & \texttt{OK} & Select smaller candidate slot\\
Malformed field combination & \texttt{INVALID\_REQUEST} & Path 255; score/quality,
state/profile zero\\
\bottomrule
\end{tabularx}
\end{table}

\section{Dwell and Failure Sequences}

The dwell mechanism is evaluated on valid accepted decisions, including valid
no-path decisions. An invalid request does not create a controller decision. The
illustrative sequence in Table~\ref{tab:worked-dwell-sequence} assumes that $\Pi_0$
is active and that the three-decision dwell has just been reset.

\begin{table}[H]
\centering
\caption{Illustrative minimum-dwell sequence. Requested profiles are hypothetical
policy-ROM outputs; the table demonstrates controller semantics only.}
\label{tab:worked-dwell-sequence}
\begin{tabularx}{\textwidth}{@{}c c c c X@{}}
\toprule
\textbf{Decision} & \textbf{Request valid} & \textbf{Requested} &
\textbf{Applied} & \textbf{Reason}\\
\midrule
1 & Yes & $\Pi_2$ & $\Pi_0$ & Minimum dwell not satisfied\\
2 & Yes, no path & $\Pi_2$ & $\Pi_0$ & Valid controller decision advances dwell\\
3 & No & --- & --- & Invalid request rejected; dwell unchanged\\
4 & Yes & $\Pi_2$ & $\Pi_2$ & Third valid dwell decision admits change\\
5 & Yes & $\Pi_1$ & $\Pi_2$ & Dwell reset after the admitted change\\
\bottomrule
\end{tabularx}
\end{table}

This sequence explains why no-path rows are included in the 81 adaptive decisions
while the three invalid physical records are not. It also shows why profile
transition counts cannot be reconstructed merely by counting changes in raw policy
ROM requests; the applied profile is a stateful result of policy and dwell.

\section{Verification Coverage Matrix}

\begin{landscape}
\begin{table}[p]
\centering
\small
\caption{Verification levels, exact covered quantities, and the defect class each
level can detect.}
\label{tab:verification-coverage-matrix}
\begin{tabularx}{\linewidth}{@{}p{3.0cm}p{4.6cm}p{3.8cm}X@{}}
\toprule
\textbf{Level} & \textbf{Registered coverage} & \textbf{Comparison rule} &
\textbf{Primary defect class}\\
\midrule
Software regression & 74 tests & Exact values and declared exceptions & Reward,
quantization, packing, path-selection logic\\
Threshold boundaries & 36 below/equal/above cases & Exact bin and state ID &
Off-by-one and real-to-code boundary errors\\
State/policy exhaustiveness & 256 state addresses and all four actions & Exact ROM
word/action & Address order, missing profile, export corruption\\
Candidate evaluator RTL & 519 of 519 vectors & Exact winner, score, and status &
Divider, saturation, range, Tchebycheff, tie-break defects\\
Integrated policy shell & 101 of 101 expected results & Exact output vector &
Pipeline alignment, H0--H4 mode selection, dwell interaction\\
Integrated cycle tracing & 2,125 cycle vectors; 530 valid decisions & Exact
accepted-edge-to-done distance & Busy/reset/wait timing and counter boundary\\
Independent implementation & Five synthesis/map/place-route/bitstream flows &
100~MHz constraint and report extraction & Configuration-specific mapping/timing defects\\
Physical replay & 420 records; 3,780 checked fields & Exact integer equality; no
tolerance & Board transport, formatter, parser, bitstream, and hardware/reference mismatch\\
\bottomrule
\end{tabularx}
\end{table}
\end{landscape}

Exactness at one level does not replace another. A Python test cannot establish
post-route timing; a timing report cannot establish UART parsing; physical
bit-exactness cannot establish route-quality superiority. The ladder therefore
supports several bounded claims rather than one undifferentiated ``validated''
label.

\section{Result-Field Accounting}

There are 84 physical records per mode, nine compared fields per record, and five
modes:
\begin{equation}
84\times9=756\ \text{fields/mode},\qquad
5\times756=3780\ \text{fields total}.
\end{equation}
The complete corpus has 76 OK, five no-path, and three invalid records per mode.
Adaptive profile counts use $84-3=81$ decisions because no-path is a valid
controller result. Held-out route-quality pairs use 58 rows because only jointly OK
rows have a selected-route quality code. These three denominators answer different
questions and should not be substituted for one another.

\section{Interpretation Checklist}

The worked cases support the following reading rules:
\begin{enumerate}
  \item A route-quality code is dimensionless UNORM16 and larger is better; it is
  not stored-key coherence.
  \item A smaller Tchebycheff score is better, but scores produced under different
  active ratios are not a common physical utility scale.
  \item A no-path result is a valid evaluated request with no feasible candidate;
  an invalid result is a malformed request rejected by the interface contract.
  \item H4's two-cycle increment over H3 is controller setup; H0's roughly
  200-cycle advantage comes from bypassing the iterative multi-objective engine.
  \item A mean reward near 4.67 is a dimensionless sum under the stated equation.
  \item Confidence intervals crossing zero and unadjusted exploratory sign tests do
  not establish route-quality superiority.
\end{enumerate}
